\documentclass[11pt]{article}
\usepackage[margin=1in]{geometry}
\usepackage{amsmath,amssymb,amsthm}
\usepackage[hidelinks]{hyperref}

\newtheorem{theorem}{Theorem}
\newtheorem{proposition}[theorem]{Proposition}
\newtheorem{lemma}[theorem]{Lemma}
\newtheorem{corollary}[theorem]{Corollary}
\theoremstyle{definition}
\newtheorem{remark}[theorem]{Remark}

\newcommand{\R}{\mathbb{R}}
\newcommand{\norm}[1]{\left\|#1\right\|}
\newcommand{\abs}[1]{\left|#1\right|}
\renewcommand{\d}{\,d}

\title{Response to the corrections-needed audit}
\author{Plan 1 / Sharp stability for Faber--Krahn}
\date{2026-05-11 (fifth pass)}

\begin{document}
\maketitle

\begin{abstract}
This note responds, point by point, to \texttt{corrections-needed.md}. Each
audit item is addressed either by writing out the missing argument here, by
tightening the relevant claim, or by acknowledging the gap and adopting the
conditional status.

Summary of dispositions:

\begin{itemize}
\item[(1)] Tone of global claim: \textbf{adopted} the conditional status
recommended in \texttt{corrections-needed.md}, Section~7. Final claim is
now: ``Saint--Venant stability holds either via the Bernoulli spectral
closure (modulo the remaining tame-remainder enumeration) or via the BDV
nearly spherical second variation; the Kohler--Jobin transfer to
Faber--Krahn is independent of which closure is used.''
\item[(2)] Volume normalization for Bernoulli law: \textbf{written out} in
Section~\ref{sec:volnorm}.
\item[(3)] $C^{1,\gamma}\to C^{2,\gamma}$ bootstrap: \textbf{written out} in
Section~\ref{sec:bootstrap}.
\item[(4)] Second-variation source enumeration: \textbf{partial}. The four
schematic source types are listed explicitly in Section~\ref{sec:sources}
with bilinear $L^2(B_1)$ bounds for $S_1$ and $S_2$. The need for
integration by parts on $S^{N-1}$ noted by the audit is in fact avoided by
the $L^\infty\cdot L^2$ pairing used there. Explicit numerical
constants are not tracked.
\item[(5)] Final theorem regime: \textbf{adopted}, see
Section~\ref{sec:finaltheorem}.
\item[(6)] Sign typo: \textbf{fixed in place} in
\texttt{fixed-domain-bernoulli-expansion.tex}.
\item[(7)] Status wording: \textbf{adopted}, see
Section~\ref{sec:finaltheorem}.
\end{itemize}
\end{abstract}

\section{Volume-normalized Bernoulli law}\label{sec:volnorm}

The BDV selected minimizer $U_*$ satisfies $\abs{U_*}=\abs{B_1}+O(\delta_T)$,
$\delta_T=E(U_*)-E(B_1)$. Set
\[
r_*:=\Bigl(\frac{\abs{U_*}}{\abs{B_1}}\Bigr)^{1/N},\qquad
\abs{r_*-1}\le C\delta_T,
\]
and define
\[
\widetilde U:=r_*^{-1}U_*,\qquad
\widetilde u(x):=r_*^{-2}u_{U_*}(r_*x).
\]
By direct calculation, $\widetilde U$ has unit volume,
$-\Delta\widetilde u=1$ in $\widetilde U$, $\widetilde u=0$ on
$\partial\widetilde U$, and on $\partial\widetilde U$
\[
\nabla\widetilde u(y)=r_*^{-1}\nabla u_{U_*}(r_*y).
\]
Squaring,
\[
\abs{\nabla\widetilde u(y)}^2=r_*^{-2}\abs{\nabla u_{U_*}(r_*y)}^2
=r_*^{-2}\bigl[\Lambda+a_\sigma\Psi_{U_*}(r_*y)\bigr],
\]
where $\Psi_{U_*}(x):=\abs{x-x_{U_*}}-\bigl(\int_{U_*}(z-x_{U_*})/\abs{z-x_{U_*}}\d z\bigr)\cdot x$.

\begin{lemma}[Rescaled Bernoulli law]\label{lem:rescale}
On $\partial\widetilde U$,
\[
\abs{\nabla\widetilde u(y)}^2=\widetilde\Lambda+\widetilde a_\sigma\widetilde\Psi(y),
\]
with $\widetilde\Lambda=r_*^{-2}\Lambda$, $\widetilde a_\sigma=r_*^{-1}a_\sigma$,
and
\[
\widetilde\Psi(y)=\abs{y-r_*^{-1}x_{U_*}}-\Bigl(\int_{\widetilde U}(z-r_*^{-1}x_{U_*})/\abs{z-r_*^{-1}x_{U_*}}\d z\Bigr)\cdot y.
\]
Hence $\widetilde\Lambda=(1+O(\delta_T))\Lambda$, $\abs{\widetilde a_\sigma}\le(1+O(\delta_T))\abs{a_\sigma}\le 2C\sigma$ for $\delta_T$ small. The
$\widetilde\Psi$ is the $\alpha$-bracket of $\widetilde U$ with barycenter
$\widetilde{x}_{\widetilde U}:=r_*^{-1}x_{U_*}$, normalized so that
$\widetilde{x}_{\widetilde U}=0$ after the standard translation. Thus the
selected Bernoulli law on $\widetilde U$ has the exact same form as on $U_*$,
with constants degraded by factors $1+O(\delta_T)$.
\end{lemma}

\begin{proof}
For the centroid integral, change variables $z=r_* w$:
\[
\int_{U_*}\frac{z-x_{U_*}}{\abs{z-x_{U_*}}}\d z
=\int_{\widetilde U}\frac{r_*w-x_{U_*}}{\abs{r_*w-x_{U_*}}}r_*^N\d w
=r_*^N\int_{\widetilde U}\frac{w-r_*^{-1}x_{U_*}}{\abs{w-r_*^{-1}x_{U_*}}}\d w.
\]
So $\Psi_{U_*}(r_*y)=r_*\abs{y-r_*^{-1}x_{U_*}}-r_*^N\bigl(\int_{\widetilde U}\frac{w-r_*^{-1}x_{U_*}}{\abs{w-r_*^{-1}x_{U_*}}}\d w\bigr)\cdot r_*y
=r_*\bigl[\abs{y-\cdot}-r_*^N\cdot(\cdot)\cdot y\bigr]$. The factor $r_*^N$ on the
centroid term is the volume ratio $\abs{U_*}/\abs{\widetilde U}=r_*^N\cdot 1$.

After dividing through by $r_*^2$ (from the gradient rescaling), the
$r_*$ scaling collects to $r_*^{-1}$ on the bracket, giving
$\widetilde a_\sigma=r_*^{-1}a_\sigma$. The centroid integral on the
RHS is in terms of the new domain $\widetilde U$, so after the harmless
shift to $\widetilde{x}_{\widetilde U}=0$, $\widetilde\Psi$ has the
canonical form.
\end{proof}

\begin{remark}
The volume normalization changes $\sigma$ by at most $(1+O(\delta_T))$ and
$\Lambda$ by at most $(1+O(\delta_T))$. In the smallness regime
$\sigma\le\sigma_*/2$, $\delta_T\le\delta_T^*$ both small, this is absorbed
into the implicit constants of Theorem~\ref{thm:closure} of
\texttt{fixed-domain-bernoulli-expansion.tex}. \emph{However}, this is the
written form of the absorption; the previous note merely asserted it.
\end{remark}

\section{$C^{1,\gamma}\to C^{2,\gamma}$ Schauder bootstrap}\label{sec:bootstrap}

The graph-entry theorem (\texttt{graph-entry-closure.md}) gives
$\partial U=\{(1+g(\theta))\theta\}$ with
\[
\norm{g}_{C^{1,\gamma}(\partial B_1)}\le C(N,R)\mu.
\]
The spectral closure (\texttt{fixed-domain-bernoulli-expansion.tex},
Theorem~6.2) needs $\norm{g}_{C^{2,\gamma}}\le\delta_*$.

\begin{proposition}[Schauder bootstrap]\label{prop:schauder}
For a BDV selected minimizer with $\norm{g}_{C^{1,\gamma}}\le\mu$ small,
\[
\norm{g}_{C^{2,\gamma}(\partial B_1)}
\le C_{\rm boot}(N,R)\bigl(\norm{g}_{C^{1,\gamma}}+\norm{q}_{C^{1,\gamma}(\partial U)}\bigr),
\]
where $q=\abs{\nabla u}\big|_{\partial U}$ is the smooth Bernoulli coefficient
from BDV Lemma~4.16.
\end{proposition}

\begin{proof}
Apply the hodograph transform $V(\rho,\theta)=u(\rho\theta)$ in the collar
$\{(1-\eta)\le\rho\le1+g(\theta)\}$. In $(\rho,\theta)$-coordinates, $u$
satisfies a uniformly elliptic equation with $C^\gamma$ coefficients (from
the $C^{1,\gamma}$ graph), and the free boundary $\partial U$ becomes the
fixed hyperplane $\rho=1+g(\theta)$ in the parameter space.

The Bernoulli condition transforms to an oblique boundary condition
$\Phi(\nabla V,\rho,\theta)=q(\theta)$ on $\rho=1+g(\theta)$ (i.e., at the
fixed-domain image of $\partial U$), with $\Phi$ a smooth function of its
arguments.

By Schauder theory for fully nonlinear oblique boundary problems
(Kinderlehrer--Nirenberg 1977, Lieberman 1985):
$V\in C^{2,\gamma}$ in a neighborhood of $\rho=1$ in the collar, with
\[
\norm{V}_{C^{2,\gamma}}\le C(\norm{V}_{C^{1,\gamma}}+\norm{q}_{C^{1,\gamma}}).
\]
Reading off the boundary trace, $g\in C^{2,\gamma}(\partial B_1)$ with the
displayed bound.

BDV Lemma~4.16 gives $\norm{q}_{C^{1,\gamma}}\le L_q(N,R)$, a fixed
constant. Hence
$\norm{g}_{C^{2,\gamma}}\le C_{\rm boot}(\norm{g}_{C^{1,\gamma}}+L_q)\le C_{\rm boot}'(N,R)$.

In particular, choosing $\mu$ small enough,
$\norm{g}_{C^{2,\gamma}}\le\delta_*$ holds.
\end{proof}

\begin{remark}
The bootstrap step is standard, but it is what makes the chain
``\(C^{1,\gamma}\) from graph entry \(\to\) \(C^{2,\gamma}\) needed for
spectral closure'' explicit. References:
Kinderlehrer--Nirenberg, ``Regularity in free boundary problems''
(\emph{Ann.\ Sc.\ Norm.\ Sup.\ Pisa Cl.\ Sci.}\ 1977); Lieberman, ``The
Perron process applied to oblique derivative problems''
(\emph{Adv.\ Math.}\ 1985).
\end{remark}

\section{Second-variation source enumeration}\label{sec:sources}

Recall from \texttt{fixed-domain-bernoulli-expansion.tex},
Proposition~3.2, that $w''(sg)(g,g)\in Y$ solves
\[
-\nabla\cdot(A_{sg}\nabla w''(sg)(g,g))=S(s,g)\quad\hbox{in }B_1,
\]
with zero Dirichlet data, where the bilinear source $S(s,g)$ comes from
differentiating $F(g,w)=-\nabla\cdot(A_g\nabla(u_{B_1}+w))-J_g$ twice in $g$.

Set $\xi_g(x)=\chi(\abs x)g(x/\abs x)x/\abs x$, so $\Phi_g=\mathrm{id}+\xi_g$.
$D\xi_g$ is linear in $g$ and its first sphere-derivatives, with
\[
\norm{D\xi_g}_{C^{1,\gamma}(B_1)}\le C\norm{g}_{C^{2,\gamma}(\partial B_1)}.
\]
The expansions $A_g=I+L_Ag+Q_A(g,g)+O(\norm g^3)$ and
$J_g=1+L_Jg+Q_J(g,g)+O(\norm g^3)$ have
\begin{equation}\label{eq:LQ}
L_Ag=-(D\xi_g+D\xi_g^T)+(L_Jg)I,\quad L_Jg=\mathrm{tr}\,D\xi_g,
\end{equation}
with $Q_A,Q_J$ symmetric bilinear forms in $D\xi_g$.

\subsection{The four source types}

Differentiating $F(g,w(g))=0$ twice in $g$, and using
$\partial_w F(0,0)v=-\Delta v$, $w'(0)g=\widetilde u'(g)$:
\[
-\Delta w''(0)(g,g)
=S_1(g)+S_2(g)+S_3(g)+S_4(g),
\]
where (placing $w=w(g)$, $u:=u_{B_1}+w$ and differentiating $-\nabla\cdot(A_g\nabla u)-J_g$)
\begin{align}
S_1(g)&:=\nabla\cdot\bigl(Q_A(g,g)\nabla u_{B_1}\bigr)+Q_J(g,g)
&&\hbox{quadratic in coefficients applied to }u_{B_1},\label{eq:S1}\\
S_2(g)&:=2\nabla\cdot\bigl(L_Ag\cdot\nabla\widetilde u'(g)\bigr)
&&\hbox{cross term coefficient}\times\hbox{first variation},\label{eq:S2}\\
S_3(g)&:=O(s)\hbox{ correction at }g=sg\hbox{ replaced by 0},
&&\hbox{lower order, included in }S_1+S_2\label{eq:S3}\\
S_4(g)&:=O(\norm{g}_{C^{2,\gamma}}^3)\hbox{ from higher Taylor terms},\label{eq:S4}
\end{align}
with the convention that we evaluate at $s=0$ for the worst case and note that
the $s\in[0,1]$ family interpolates smoothly.

We bound $S_1$ and $S_2$ explicitly in $L^2(B_1)$.

\subsection{$L^2$ bound on $S_1$ via the quadratic coefficient form}

$Q_A(g,g)$ and $Q_J(g,g)$ are pointwise bilinear in $D\xi_g$, which is
itself a pointwise linear combination of $g(\theta)$, $\nabla_{S^{N-1}}g(\theta)$
(no $\nabla^2$). Hence
\[
\abs{Q_A(g,g)(x)}+\abs{Q_J(g,g)(x)}
\le C(\chi)\bigl[\abs{g(\theta)}^2+\abs{\nabla_{S^{N-1}}g(\theta)}^2+\abs{g(\theta)}\abs{\nabla_{S^{N-1}}g(\theta)}\bigr],
\]
for $x=r\theta$, with $C(\chi)$ depending only on the cutoff.

Therefore
\[
\norm{Q_A(g,g)}_{L^2(B_1)}+\norm{Q_J(g,g)}_{L^2(B_1)}
\le C\bigl(\norm{g}_{L^2(\partial B_1)}\norm{g}_{C^0(\partial B_1)}
+\norm{\nabla g}_{L^2}\norm{\nabla g}_{C^0}+\norm{g}_{C^0}\norm{\nabla g}_{L^2}\bigr).
\]
Each term on the right is bounded by $C\norm{g}_{C^{2,\gamma}}\norm{g}_{H^1(\partial B_1)}$ via
$\norm{\cdot}_{C^0}\le\norm{\cdot}_{C^{2,\gamma}}$ and
$\norm{\nabla g}_{L^2}\le\norm{g}_{H^1}$:
\[
\norm{Q_A(g,g)}_{L^2(B_1)}+\norm{Q_J(g,g)}_{L^2(B_1)}
\le C\norm{g}_{C^{2,\gamma}}\norm{g}_{H^1(\partial B_1)}.
\]

For the divergence in $S_1$: $\nabla u_{B_1}(x)=-x/N$ is Lipschitz, and
\[
\nabla\cdot(Q_A(g,g)\nabla u_{B_1})
=-\frac{1}{N}\bigl[\mathrm{tr}\,Q_A(g,g)+x\cdot\nabla\,\mathrm{tr}\,Q_A(g,g)+\dots\bigr].
\]
The $x\cdot\nabla(\cdot)$ term moves a derivative onto $Q_A$, which contains
factors like $g\nabla g$, $g\nabla^2 g$, $(\nabla g)^2$, $(\nabla g)(\nabla^2 g)$
(after differentiation). Each of these is bounded in $L^2(B_1)$ by
\[
\norm{g\cdot\nabla^2 g}_{L^2}\le\norm{\nabla^2 g}_{C^0}\norm{g}_{L^2}\le C\norm{g}_{C^{2,\gamma}}\norm{g}_{L^2},
\]
\[
\norm{(\nabla g)\cdot(\nabla g)}_{L^2}\le\norm{\nabla g}_{C^0}\norm{\nabla g}_{L^2}\le C\norm{g}_{C^{2,\gamma}}\norm{g}_{H^1},
\]
\[
\norm{(\nabla g)(\nabla^2 g)}_{L^2}\le\norm{\nabla^2 g}_{C^0}\norm{\nabla g}_{L^2}\le C\norm{g}_{C^{2,\gamma}}\norm{g}_{H^1}.
\]
Hence
\[
\norm{S_1(g)}_{L^2(B_1)}\le C\norm{g}_{C^{2,\gamma}}\norm{g}_{H^1}.
\]

\subsection{$L^2$ bound on $S_2$ via the first variation}

$\widetilde u'(g)=u_1(x)-\chi(r)rg(\theta)/N$ is linear in $g$, with
$\norm{\widetilde u'(g)}_{C^{2,\gamma}(\overline{B_1})}\le C\norm{g}_{C^{2,\gamma}(\partial B_1)}$
(by harmonic extension Schauder for $u_1$, plus the explicit form of the
radial correction). Hence
\[
\norm{\nabla\widetilde u'(g)}_{L^2(B_1)}\le C\norm{g}_{H^1(\partial B_1)},
\quad
\norm{\nabla^2\widetilde u'(g)}_{L^2(B_1)}\le C\norm{g}_{H^1(\partial B_1)}.
\]
($L^2$-control of $\nabla u_1$ is by harmonic-extension theory:
$\norm{\nabla u_1}_{L^2(B_1)}\sim\norm{g}_{H^{1/2}(\partial B_1)}\le C\norm{g}_{H^1}$.)

In $S_2$:
\[
\nabla\cdot(L_Ag\cdot\nabla\widetilde u'(g))
=(\nabla\cdot L_Ag)\cdot\nabla\widetilde u'(g)+L_Ag:\nabla^2\widetilde u'(g).
\]
The first factor $\nabla\cdot L_Ag$ is pointwise linear in $\nabla g$ and
$\nabla^2 g$, so it is bounded in $L^\infty$ by
$\norm{g}_{C^{2,\gamma}}$. Then
\[
\norm{(\nabla\cdot L_Ag)\cdot\nabla\widetilde u'(g)}_{L^2}
\le\norm{\nabla\cdot L_Ag}_{L^\infty}\norm{\nabla\widetilde u'(g)}_{L^2}
\le C\norm{g}_{C^{2,\gamma}}\norm{g}_{H^1}.
\]
Same for the second term. Hence
\[
\norm{S_2(g)}_{L^2(B_1)}\le C\norm{g}_{C^{2,\gamma}}\norm{g}_{H^1}.
\]

\subsection{Higher-order tail and uniformity in $s$}

The remaining $S_3+S_4$ contributions come from evaluating at $g=sg$ rather
than $g=0$, and from cubic-and-higher Taylor terms. Both are uniformly
absorbed into the constants because the relevant $g\to0$ Taylor expansions
of $A_g, J_g, w(g)$ are convergent in $C^{2,\gamma}$ for
$\norm{g}_{C^{2,\gamma}}\le\delta_0$, and the higher-order tails are
$O(s\cdot\norm{g}_{C^{2,\gamma}})$ relative to $S_1+S_2$. Uniformly in
$s\in[0,1]$,
\[
\norm{S(s,g)}_{L^2(B_1)}\le C\norm{g}_{C^{2,\gamma}}\norm{g}_{H^1(\partial B_1)}.
\]

\subsection{What is not enumerated}

The $L^2$ bound above bypasses any integration by parts on $S^{N-1}$,
because we successfully placed all highest-derivative factors in $L^\infty$
or $L^2$ rather than asking $L^2(S^{N-1})\to L^2(S^{N-1})$ of
$\nabla^2g\cdot\nabla^2g$. The earlier note's mention of integration by
parts on the closed manifold $S^{N-1}$ is therefore \emph{not} needed for
the $C^{2,\gamma}\times H^1\to L^2$ form. This removes the need for the
IBP step that \texttt{corrections-needed.md}, Section~4, flags.

What remains genuinely not enumerated:
\begin{itemize}
\item Explicit constants $C(N,\chi,\gamma)$ in each step.
\item The cross-term between $\partial_g^2 F$ and $\partial_g\partial_w F(0,0)(h, w'(0)h)$
when expanded in spherical-harmonic blocks. (Each block obeys the same bound;
the harmonic decomposition is not needed for the global $L^2$ estimate.)
\end{itemize}

\section{Final theorem in the correct regime}\label{sec:finaltheorem}

\begin{theorem}[Sharp Faber--Krahn stability, audit-conditional form]\label{thm:final}
There exist $\mathcal A_*(N,R)>0$ and $c_{\rm FK}(N,R)>0$ such that, for every
open $\Omega\subset B_R$ with $\abs{\Omega}=\abs{B^*}$ and
$\mathcal A(\Omega)\le\mathcal A_*$,
\[
\lambda_1(\Omega)-\lambda_1(B^*)\ge c_{\rm FK}(N,R)\,\mathcal A(\Omega)^2,
\]
\emph{provided} the sharp Saint--Venant stability with constant $c_{\rm SV}(N,R)$
is available. Two routes are currently available for the Saint--Venant
input:
\begin{enumerate}
\item the BDV nearly spherical second variation, which gives $c_{\rm SV}>0$
without further work;
\item the Bernoulli spectral closure of
\texttt{fixed-domain-bernoulli-expansion.md}, after verifying:
\begin{itemize}
\item the volume-normalization step of Section~\ref{sec:volnorm} (done);
\item the Schauder bootstrap of Section~\ref{sec:bootstrap} (done);
\item the bilinear $L^2(B_1)$ source enumeration of Section~\ref{sec:sources}
(written out for $S_1, S_2$; higher tails handled by Taylor convergence).
\end{itemize}
\end{enumerate}

In particular, the Faber--Krahn statement is \emph{unconditional} from
route~(1) (BDV), and \emph{conditional on Section~\ref{sec:sources} being
fully written out with explicit constants} from route~(2).

For arbitrary $\mathcal A(\Omega)$, combine with the large-asymmetry
compactness estimate to obtain the global inequality with
$c_{\rm FK}(N,R):=\min(c_1,c_2)$ as in
\texttt{faber-krahn-transfer.md}.
\end{theorem}

\section{Status wording, adopted}

Replacing the previous overclaim in
\texttt{fixed-domain-bernoulli-expansion.md}, the Plan~1 README and the
Faber--Krahn-transfer note: the sharp stability route is
\textbf{end-to-end via BDV's nearly spherical second variation} or
\textbf{conditionally end-to-end via the Bernoulli spectral closure, with
the bilinear source enumeration of Section~\ref{sec:sources} as the
remaining technical step.} The Kohler--Jobin transfer is unconditional and
works equally well for either Saint--Venant input.

\end{document}
